\chapter{Симметрия, вырождение и лестницы углового момента}
\providecommand{\ket}[1]{\left\lvert#1\right\rangle}

\section{Симметрии и вырожденные уровни энергии}

Симметрия — это операция, которая не изменяет физическое описание системы. Она не изменяет гамильтониан, уравнение Schr\"odinger или возможные энергии. Если эту операцию представляет унитарный оператор \(S\), то инвариантность означает
\begin{equation}
S H S^\dagger=H,
\end{equation}
или, что эквивалентно,
\begin{equation}
[S,H]=0.
\end{equation}

Перенос — один из примеров. Если переместить атом из одного места в другое, его энергия не изменится. Это верно для любого атомного состояния, а не только для основного. Другой пример — вращение, и именно оно будет здесь главным предметом.

Перестановка двух тождественных частиц тоже является симметрией. Назвать два электрона первым и вторым, или Harry и Fred, не значит сделать их различными. Перестановка меток не может изменить свойства атома.

У кристалла есть дополнительный вид трансляционной симметрии. Перемещение всего кристалла в пространстве — это обычный перенос. Однако для бесконечного кристалла перенос на один период решётки тоже возвращает ту же физическую конфигурацию. Эквивалентно, узлы решётки можно перенумеровать, сдвинув индекс каждого узла на единицу. Границы нарушают это точное утверждение для конечного кристалла, но вдали от концов очень длинного кристалла такая симметрия становится хорошей симметрией.

Однородность пространства---тот факт, что любое положение эквивалентно любому другому---и есть симметрия относительно переносов. Изотропия пространства---тот факт, что любое направление эквивалентно любому другому---и есть вращательная симметрия.

Вырождение связано с симметрией, но не тождественно ей. Энергетический уровень вырожден, когда более чем одно состояние имеет в точности эту энергию. Для общего гамильтониана с дискретным спектром точное совпадение двух не связанных между собой уровней было бы поразительной случайностью. Совпадение с точностью до одной части на \(10^{15}\) всё же не является точным равенством. Точные вырождения обычно указывают на симметрию.

Симметрия иногда влечёт вырождение, но не всегда. Первый пример показывает и почему это возможно, и почему одной симметрии самой по себе может быть недостаточно.

\section{Частица на окружности}

Рассмотрим бусину, вынужденную двигаться по круговой проволоке. Её положение можно было бы описывать декартовыми координатами \(x\) и \(y\), но достаточно одной угловой координаты:
\begin{equation}
\psi=\psi(\theta).
\end{equation}
Плотность вероятности положения на окружности равна \(\psi^*(\theta)\psi(\theta)\).

Активный поворот против часовой стрелки на малый угол \(\epsilon\) переводит
\begin{equation}
\psi(\theta)\longrightarrow\psi(\theta-\epsilon).
\end{equation}
Следовательно,
\begin{align}
\delta\psi(\theta)
&=\psi(\theta-\epsilon)-\psi(\theta)\\
&=-\epsilon\frac{d\psi}{d\theta}
+\mathcal O(\epsilon^2).
\end{align}

Оператор
\begin{equation}
L_z=-i\hbar\frac{d}{d\theta}
\end{equation}
эрмитов на периодических волновых функциях. Он является угловым аналогом
\begin{equation}
p_x=-i\hbar\frac{d}{dx}.
\end{equation}
Множитель \(i\) существен для эрмитовости. Через \(L_z\)
\begin{equation}
\delta\psi
=-\frac{i\epsilon}{\hbar}L_z\psi
+\mathcal O(\epsilon^2),
\end{equation}
поэтому оператор бесконечно малого поворота имеет вид
\begin{equation}
R(\epsilon)
=1-\frac{i\epsilon}{\hbar}L_z
+\mathcal O(\epsilon^2).
\end{equation}
Тем самым \(L_z\) отождествляется с генератором вращений вокруг оси, перпендикулярной окружности.

\begin{classroomqa}
\lecturetimestamp{00:11:19}
\audiencequestion{Не должен ли один из множителей с \(i\) иметь противоположный знак?}
\lecturerresponse{Проверка знаков даёт приведённое выше соотношение. При соглашении \(\psi(\theta)\to\psi(\theta-\epsilon)\) изменение равно \(-\epsilon\,d\psi/d\theta\), что совпадает с \(-i\epsilon L_z\psi/\hbar\).}
\end{classroomqa}

Теперь найдём собственные состояния углового момента:
\begin{equation}
L_z\psi=M\psi.
\end{equation}
Подстановка дифференциального оператора даёт
\begin{equation}
-i\hbar\frac{d\psi}{d\theta}=M\psi,
\end{equation}
решения которого имеют вид
\begin{equation}
\psi(\theta)=C\exp\left(\frac{iM\theta}{\hbar}\right).
\end{equation}

Есть дополнительное условие, которого не возникает при движении по бесконечной прямой. Точки \(\theta=0\) и \(\theta=2\pi\) — это одна и та же точка, поэтому однозначная орбитальная волновая функция должна удовлетворять условию
\begin{equation}
\psi(\theta+2\pi)=\psi(\theta).
\end{equation}
Для приведённой выше собственной функции это требует
\begin{equation}
\exp\left(\frac{2\pi iM}{\hbar}\right)=1.
\end{equation}
Следовательно,
\begin{equation}
M=m\hbar,
\qquad
m\in\mathbb Z.
\end{equation}
Поэтому возможные орбитальные угловые моменты на окружности таковы:
\begin{equation}
\ldots,-2\hbar,-\hbar,0,\hbar,2\hbar,\ldots.
\end{equation}

Должно ли состояние с угловым моментом \(m\hbar\) иметь ту же энергию, что и состояние с угловым моментом \(-m\hbar\)? Положительные и отрицательные \(m\) соответствуют противоположным направлениям обращения. Вопрос в том, требует ли одна лишь вращательная симметрия, чтобы эти два движения имели одинаковую энергию.

\begin{classroomqa}
\lecturetimestamp{00:22:51}
\audiencequestion{Означает ли использование \(\psi(\theta-\epsilon)\), что частица движется по часовой стрелке?}
\lecturerresponse{Нет. Это соглашение лишь фиксирует, какое направление называется положительным вращением. Оно определяет знаки угловой координаты и углового момента, но не задаёт фактическое направление, в котором движется частица.}
\end{classroomqa}

Магнитное поле, перпендикулярное плоскости окружности, даёт контрпример. Оно может смещать энергии состояний с положительным и отрицательным угловым моментом в противоположные стороны и тем самым расщеплять эту пару. Знак смещения зависит от заряда и соглашений об ориентации, но важен именно сам факт расщепления.

Такое поле всё ещё сохраняет вращения вокруг перпендикулярной оси. Поэтому оно показывает, что вращательная симметрия сама по себе не требует
\begin{equation}
E(m)=E(-m).
\end{equation}
При отсутствии расщепления \(m=1\) и \(m=-1\) могут образовывать вырожденную пару, как и \(m=2\) и \(m=-2\). Состояние \(m=0\) не имеет отличного от него партнёра такого рода.

\section{Зеркальная симметрия и пара \texorpdfstring{\(m,-m\)}{m, -m}}

Добавим дискретную зеркальную симметрию. Пусть \(\mathcal M\) обозначает отражение относительно диаметра окружности:
\begin{equation}
(\mathcal M\psi)(\theta)=\psi(-\theta).
\end{equation}
Оно меняет направление обращения на противоположное:
\begin{equation}
\mathcal M\psi_m=\psi_{-m}.
\end{equation}

Если отражение является симметрией, то
\begin{equation}
[\mathcal M,H]=0.
\end{equation}
Предположим,
\begin{equation}
H\ket{m}=E_m\ket{m}.
\end{equation}
Применяя \(\mathcal M\), получаем
\begin{align}
H\mathcal M\ket{m}
&=\mathcal M H\ket{m}\\
&=E_m\mathcal M\ket{m}.
\end{align}
Поскольку \(\mathcal M\ket{m}\) — это состояние с угловым моментом \(-m\hbar\),
\begin{equation}
E_{-m}=E_m.
\end{equation}
Вращательная симметрия сохраняет угловой момент; зеркальная симметрия связывает два противоположных значения. Вместе они требуют, чтобы эта пара была вырожденной.

\begin{classroomqa}
\lecturetimestamp{00:29:54}
\audiencequestion{Поскольку моды \(m\) и \(-m\) имеют одну и ту же плотность вероятности, показывает ли это уже, что их энергии равны?}
\lecturerresponse{Нет. Волновые функции \(e^{im\theta}\) и \(e^{-im\theta}\) различны, хотя их модули в квадрате совпадают. Уравнение на собственные значения углового момента включает волновую функцию, а не только \(\psi^*\psi\). Квантовая механика требует амплитуд, а не одних лишь вероятностей.}
\end{classroomqa}

Магнитное поле сохраняет вращения вокруг своей оси, но нарушает использованную выше зеркальную симметрию.

\begin{classroomqa}
\lecturetimestamp{00:33:05}
\audiencequestion{Меняет ли зеркальная симметрия направление магнитного поля на противоположное?}
\lecturerresponse{Да. Магнитное поле — аксиальный вектор. Если отразить поле, направленное в плоскость рисунка, получится поле, направленное из плоскости рисунка. Эквивалентно, токовая петля и её зеркальное изображение несут ток в противоположных направлениях и создают противоположно направленные поля.}
\end{classroomqa}

Некоммутативность вращения и отражения можно увидеть непосредственно. На собственной функции углового момента,
\begin{align}
\mathcal M L_z\psi_m
&=m\hbar\,\psi_{-m},\\
L_z\mathcal M\psi_m
&=-m\hbar\,\psi_{-m}.
\end{align}
Следовательно,
\begin{equation}
[\mathcal M,L_z]\psi_m
=2m\hbar\,\psi_{-m},
\end{equation}
что ненулево при \(m\neq0\). Сначала повернуть, а затем отразить — не то же самое, что сначала отразить, а затем повернуть.

Когда оба преобразования являются симметриями, оба оператора коммутируют с \(H\), но друг с другом они не коммутируют. Именно такая структура отвечает за вырождение.

\begin{classroomqa}
\lecturetimestamp{00:40:28}
\audiencequestion{Порождает ли вырождения любая некоммутирующая гейзенбергова пара, такая как \(x\) и \(p\)?}
\lecturerresponse{Вообще говоря, нет. Операторы должны также быть симметриями: они должны коммутировать с гамильтонианом. Координата и импульс обычно не сохраняются одновременно. Если бы оба коммутировали с \(H\), то их некоммутативность влекла бы вырождение.}
\end{classroomqa}

\section{Коммутаторные алгебры симметрий}

Предположим, что \(A\) и \(B\) — эрмитовы генераторы двух непрерывных симметрий:
\begin{equation}
[A,H]=0,
\qquad
[B,H]=0.
\end{equation}
Ничто в этих уравнениях не требует, чтобы \(A\) и \(B\) коммутировали друг с другом.

Коммутатор \([A,B]\) антиэрмитов, поэтому определим эрмитов оператор
\begin{equation}
C=i[A,B].
\end{equation}
Численные множители и общий знак здесь несущественны. Важно то, что \(C\) тоже коммутирует с гамильтонианом. Действительно,
\begin{equation}
[AB,H]
=A[B,H]+[A,H]B=0,
\end{equation}
и аналогично
\begin{equation}
[BA,H]=0.
\end{equation}
Следовательно,
\begin{equation}
[C,H]
=i\bigl([AB,H]-[BA,H]\bigr)=0.
\end{equation}

Новый генератор \(C\) может быть независимым от \(A\) и \(B\), а может оказаться лишь линейной комбинацией уже известных генераторов. Если он новый, добавьте его в список и образуйте дальнейшие коммутаторы. Повторение этого процесса может продолжаться бесконечно, но часто он замыкается после конечного числа генераторов: каждый следующий коммутатор является линейной комбинацией уже имеющихся. Такая замкнутая система называется коммутаторной алгеброй, или алгеброй Ли.

Конечные преобразования симметрии образуют группу. Их бесконечно малые генераторы образуют соответствующую алгебру Ли.

Вырождения важны, потому что они организуют спектры. Если два состояния имеют одинаковую энергию, одно не может перейти в другое, испустив фотон, несущий разность их энергий. В атомной спектроскопии расположение уровней и их вырождение определяют возможные спектральные линии. Та же логика применима к возбуждённым адронам и их переходам.

\begin{classroomqa}
\lecturetimestamp{00:49:07}
\audiencequestion{Являются ли сами эти коммутаторы элементами группы симметрии?}
\lecturerresponse{Обсуждаемые операторы являются генераторами. Генераторы и их коммутаторы образуют коммутаторную алгебру, или алгебру Ли. Конечные операции симметрии, построенные из генераторов, являются элементами группы.}
\end{classroomqa}

Если все генераторы симметрии коммутируют друг с другом, симметрия называется абелевой. Если некоторые их коммутаторы ненулевые, она неабелева.

\begin{classroomqa}
\lecturetimestamp{00:50:25}
\audiencequestion{Каково фундаментальное определение генератора?}
\lecturerresponse{Возьмите унитарное преобразование симметрии, очень близкое к тождественному, и запишите его в виде \(1-i\epsilon G/\hbar+\cdots\). Эрмитов оператор \(G\) является его генератором. Повторяя соответствующее бесконечно малое преобразование, строят конечное.}
\end{classroomqa}

\begin{classroomqa}
\lecturetimestamp{00:52:08}
\audiencequestion{Что понимается под алгеброй генераторов симметрии?}
\lecturerresponse{Это набор генераторов, которые можно складывать, вычитать, умножать на обычные числа и брать их коммутаторы. В этой алгебре коммутатор играет роль умножения, а замкнутость означает, что эти операции не порождают генераторов вне этого набора.}
\end{classroomqa}

\section{Повороты вокруг разных осей}

Повороты дают важнейший пример неабелевой симметрии. Поворот вокруг оси \(x\), за которым следует поворот вокруг оси \(y\), вообще говоря, не совпадает с выполнением их в обратном порядке. Их несовпадение само связано с поворотом вокруг третьей оси.

\begin{classroomqa}
\lecturetimestamp{00:54:22}
\audiencequestion{Коммутируют ли бесконечно малые повороты, хотя конечные повороты не коммутируют?}
\lecturerresponse{Они различаются во втором порядке по двум малым углам поворота. Этого достаточно: после вынесения этих малых параметров сами генераторы имеют ненулевой коммутатор.}
\end{classroomqa}

Прежде чем вычислять квантовые коммутаторы, рассмотрим соответствующую классическую картину. Возьмём круговую планетную орбиту, угловой момент которой перпендикулярен её плоскости. Поворот всей орбиты вокруг оси \(y\) даёт другую допустимую орбиту с точно такой же энергией. Однако её вектор углового момента теперь имеет другую \(x\)-компоненту.

Таким образом, вращательная инвариантность порождает конфигурации с одинаковой энергией и разными компонентами углового момента. В квантовой механике это проявляется как некоммутативность между генераторами поворотов вокруг разных осей.

\section{Угловой момент в декартовых координатах}

Сначала рассмотрим частицу, движущуюся в плоскости \(xy\), не ограничивая её движением по окружности. При малом повороте против часовой стрелки,
\begin{equation}
\delta x=-\epsilon y,
\qquad
\delta y=\epsilon x.
\end{equation}
Для активно повернутого состояния преобразованная волновая функция вычисляется в координатах, повернутых в обратную сторону:
\begin{equation}
\psi_\epsilon(x,y)
=\psi(x+\epsilon y,y-\epsilon x).
\end{equation}
Следовательно,
\begin{align}
\delta\psi
&=\epsilon\left(
y\frac{\partial}{\partial x}
-x\frac{\partial}{\partial y}
\right)\psi
+\mathcal O(\epsilon^2)\\
&=-\frac{i\epsilon}{\hbar}
\left(xp_y-yp_x\right)\psi
+\mathcal O(\epsilon^2).
\end{align}
Поэтому генератор поворотов вокруг оси \(z\) имеет вид
\begin{equation}
L_z=xp_y-yp_x.
\end{equation}
Это \(z\)-компонента знакомого классического вектора
\begin{equation}
\mathbf L=\mathbf r\times\mathbf p.
\end{equation}
Это не только механический угловой момент; в квантовой механике это оператор, порождающий повороты волновой функции.\footnote{Выбрано соглашение об активном повороте, чтобы оно согласовывалось с более ранним преобразованием \(\psi(\theta)\mapsto\psi(\theta-\epsilon)\) и с окончательным отождествлением \(L_z=xp_y-yp_x\) после промежуточной неопределённости знака в устном выводе.}

Циклическая перестановка \(x\to y\to z\to x\) даёт остальные компоненты.\footnote{Формулы для \(L_x\) и \(L_y\) восстанавливаются применением циклической подстановки \(x\to y\to z\to x\), явно предписанной на лекции, к \(L_z=xp_y-yp_x\).}
\begin{align}
L_x&=yp_z-zp_y,\\
L_y&=zp_x-xp_z,\\
L_z&=xp_y-yp_x.
\end{align}

Для системы, инвариантной относительно поворотов, эти операторы коммутируют с гамильтонианом:
\begin{equation}
[H,L_x]=[H,L_y]=[H,L_z]=0.
\end{equation}
Стандартный пример — частица в поле центральной силы. Эти уравнения представляют собой закон сохранения углового момента в квантовой форме.

\section{Алгебра Ли углового момента}

Канонические коммутационные соотношения имеют вид
\begin{align}
[x_i,x_j]&=0,\\
[p_i,p_j]&=0,\\
[x_i,p_j]&=i\hbar\,\delta_{ij}.
\end{align}
Координаты коммутируют с компонентами импульса, относящимися к другим направлениям. Так, \(x\) коммутирует с \(p_y\), например, тогда как \(x\) не коммутирует с \(p_x\).

Теперь вычислим
\begin{equation}
[L_x,L_y]
=
[yp_z-zp_y,\;zp_x-xp_z].
\end{equation}
Большинство членов обращается в нуль. Два ненулевых вклада равны
\begin{align}
[yp_z,zp_x]&=-i\hbar\,yp_x,\\
[zp_y,xp_z]&=i\hbar\,xp_y.
\end{align}
Следовательно,
\begin{align}
[L_x,L_y]
&=i\hbar\left(xp_y-yp_x\right)\\
&=i\hbar L_z.
\end{align}
Циклическая перестановка координатных меток даёт
\begin{align}
[L_x,L_y]&=i\hbar L_z,\\
[L_y,L_z]&=i\hbar L_x,\\
[L_z,L_x]&=i\hbar L_y.
\end{align}

\begin{classroomqa}
\lecturetimestamp{01:14:50}
\audiencequestion{Когда коммутатор даёт \(L_z\), нужно ли нам ещё показать, что \(L_z\) не является линейной комбинацией \(L_x\) и \(L_y\)?}
\lecturerresponse{Да. Следует установить эту независимость, а не просто предполагать её. Прямое алгебраическое доказательство на данном этапе не рассматривается.}
\end{classroomqa}

Выписанные коммутаторы замыкаются на линейной оболочке \(L_x\), \(L_y\) и \(L_z\): как только включены все три, дальнейшие коммутаторы не дают ни одного оператора вне этой оболочки. Это алгебра Ли углового момента. Вместе с
\begin{equation}
[H,L_i]=0,
\qquad i=x,y,z,
\end{equation}
эти соотношения содержат существенную информацию об энергетическом спектре, не требуя прямого решения угловой части уравнения Schr\"odinger.

Аналогичная алгебраическая стратегия используется для квантового гармонического осциллятора, где операторы рождения и уничтожения порождают его спектр. Здесь аналогичные операторы будут изменять \(z\)-компоненту углового момента, а не энергию.

Выбор оси \(z\) — лишь выбор представления. Ничто не делает \(z\) физически предпочтительнее, чем \(x\) или \(y\), но выбор одной компоненты даёт удобный базис. Запишем
\begin{equation}
L_z\ket{m}=m\hbar\ket{m}.
\end{equation}
Задача состоит в том, чтобы определить, какие значения \(m\) могут встречаться и как связаны состояния с различными \(m\).

\begin{classroomqa}
\lecturetimestamp{01:24:02}
\audiencequestion{Сравнивая эту конструкцию с гармоническим осциллятором, вы имеете в виду квантовый гармонический осциллятор? Это не разбиралось на предыдущем занятии.}
\lecturerresponse{Да, сравнение проводится с квантовым осциллятором. Настоящее рассуждение об угловом моменте самодостаточно и не требует знания осциллятора. Позднее аналогию можно читать и в обратную сторону: в осцилляторе используется сходный алгебраический метод.}
\end{classroomqa}

\section{Повышение и понижение углового момента}

Определим
\begin{equation}
L_+=L_x+iL_y,
\qquad
L_-=L_x-iL_y.
\end{equation}
Из коммутационных соотношений углового момента следует
\begin{equation}
[L_z,L_+]=\hbar L_+,
\qquad
[L_z,L_-]=-\hbar L_-.
\end{equation}

Применим первое соотношение к собственному состоянию \(L_z\):
\begin{align}
L_zL_+\ket{m}
&=\left(L_+L_z+[L_z,L_+]\right)\ket{m}\\
&=(m+1)\hbar\,L_+\ket{m}.
\end{align}
Если \(L_+\ket{m}=0\) не выполняется, это другое собственное состояние \(L_z\), теперь с собственным значением \((m+1)\hbar\). Аналогично,
\begin{equation}
L_zL_-\ket{m}
=(m-1)\hbar\,L_-\ket{m}.
\end{equation}
Тем самым \(L_+\) повышает \(m\) на единицу, а \(L_-\) понижает его на единицу.

Начиная с любого одного собственного значения, алгебра порождает последовательность, значения в которой отличаются на целые числа. Она может продолжаться неограниченно, а может обрываться, поскольку
\begin{equation}
L_+\ket{m_{\mathrm{top}}}=0
\end{equation}
на верхнем конце и
\begin{equation}
L_-\ket{m_{\mathrm{bottom}}}=0
\end{equation}
на нижнем.

\begin{classroomqa}
\lecturetimestamp{01:35:11}
\audiencequestion{Разве из аргумента с окружностью уже не следовало, что угловой момент обязан быть целым?}
\lecturerresponse{Этот аргумент доказывает целые значения орбитального углового момента, представленного однозначными волновыми функциями положения. Абстрактная алгебра углового момента даёт более слабый результат: лестница может содержать либо целые, либо полуцелые значения.}
\end{classroomqa}

Если лестница обрывается, вращательная инвариантность связывает её два конца. Поворот на \(180^\circ\) вокруг оси, перпендикулярной \(z\), переводит
\begin{equation}
L_z\longrightarrow-L_z.
\end{equation}
Следовательно, спектр должен быть симметричен относительно нуля:
\begin{equation}
m_{\mathrm{bottom}}=-m_{\mathrm{top}}.
\end{equation}
Поскольку соседние значения различаются на единицу, имеются две возможности. Лестница содержит либо целые числа,
\begin{equation}
\ldots,-2,-1,0,1,2,\ldots,
\end{equation}
либо полуцелые,
\begin{equation}
\ldots,-\frac{5}{2},-\frac{3}{2},-\frac{1}{2},
\frac{1}{2},\frac{3}{2},\frac{5}{2},\ldots,
\end{equation}
причём её концы находятся там, где оператор повышения или понижения даёт нуль.

Аргумент об однозначности волновой функции на окружности выбирает целочисленную возможность для орбитального углового момента,
\begin{equation}
\mathbf L=\mathbf r\times\mathbf p.
\end{equation}
Полуцелая возможность тоже физически реализуется и относится к спиновому угловому моменту, который будет рассмотрен отдельно.

Наконец, лестничные состояния имеют одинаковую энергию, если гамильтониан инвариантен относительно поворотов. Предположим, что
\begin{equation}
H\ket{m}=E\ket{m}.
\end{equation}
Поскольку каждая компонента \(L_i\) коммутирует с \(H\), с ним коммутируют и \(L_+\), и \(L_-\). Следовательно,
\begin{align}
H L_+\ket{m}
&=L_+H\ket{m}\\
&=E L_+\ket{m}.
\end{align}
Всякий раз, когда \(L_+\ket{m}\neq0\), поднятое состояние имеет ту же энергию, что и исходное состояние. Тот же расчёт даёт
\begin{equation}
H L_-\ket{m}=E L_-\ket{m}.
\end{equation}
Следовательно, операторы повышения и понижения порождают семейство состояний с различными собственными значениями \(L_z\), но с совершенно одинаковой энергией. Такова связь между некоммутирующими симметриями и вырождением, развивавшаяся на протяжении всей лекции.

\begin{classroomqa}
\lecturetimestamp{01:44:13}
\audiencequestion{Связаны ли эти лестничные состояния с выбором разных осей вращения?}
\lecturerresponse{Они связаны поворотами. Собственные векторы \(L_x\) или \(L_y\) являются линейными комбинациями выбранных собственных векторов \(L_z\). Поворот состояния изменяет его компоненты углового момента, оставляя его энергию неизменной.}
\end{classroomqa}

Выбор \(L_+\) и \(L_-\) удобен, потому что эти комбинации непосредственно выявляют лестничную структуру. Этот прием повсюду встречается в квантовой механике: когда алгебра коммутаторов замыкается, операторы повышения и понижения часто упорядочивают соответствующий спектр собственных значений.

\begin{classroomqa}
\lecturetimestamp{01:46:33}
\audiencequestion{Что происходит, когда \(L_+\) действует на состояние с максимальным значением, и как определяется верхний конец?}
\lecturerresponse{Он дает нуль. Это и есть определение верхнего конца обрывающейся лестницы. Аналогично, \(L_-\) дает нуль на нижнем конце. Абстрактная лестница не обязана обрываться. Для рассматриваемых далее состояний атомного углового момента обычно встречаются обрывающиеся лестницы.}
\end{classroomqa}

Следующий шаг — дальше развить эту структуру углового момента и применить ее к атомам, включая водород. Найденные здесь вращательные мультиплеты объясняют часть структуры вырождения; дальнейшую структуру еще предстоит разработать.
